\section{Limitations}
\label{sec:limitations}

We state the threats to validity plainly, in the spirit of the
pre-registration above.

\paragraph{Inside the adversarial regime, the mechanism is not the only
answer either.} We pre-registered that no strike count would hold both
axes under attack and it was false: a forgiveness counter at $k{=}5$
matches the ledger's corner at two lies per cycle
(\S\ref{sec:countersweep}). Two things follow, and only the second is a
defence. The first is that our headline counter numbers were a statement
about four chosen settings, and any claim of the form ``counters lose''
must now be read as scoped to the lifetime-strike shape and to untuned
forgiveness counters. The second is that the tuned counter reaches the
corner by barely evicting---$14.9$ surviving entries against the
ledger's $4.0$, and $16\%$ less measured resource at every budget---so
the two mechanisms are not interchangeable, they price different things.
A reader who does not pay for store size should use the counter.

The sweep has its own limit: it moves one parameter, at seven values, on
one environment family, and says nothing about a differently
\emph{shaped} heuristic. Patching a counter with decay and magnitude
grading reinvents the ledger, which is an argument about where the design
space leads and not a measurement of it.

\paragraph{Outside the adversarial regime, the mechanism may be
unnecessary.}
We state the strongest argument against this paper because our own data
makes it. Three independent tests failed to find an advantage for the
ledger where no adversary is present: pre-registered gates on trajectory
selection found it matched or beaten by a threshold counter under
i.i.d.\ noise, by a buffer frontier, and by a tuned EWMA under drift;
with a real model deciding adoption a one-line eviction counter revokes
poison sooner (\S\ref{sec:wef}); and on real software-engineering tasks
it shows no learning-curve advantage over a token-matched random control,
with the no-memory arm scoring highest of five (\S\ref{sec:swebench}).
Each has a local explanation, and we give those where they belong. But a
reader is entitled to the pattern rather than three separate mitigations:
absent an attacker, what we have demonstrated is leanness and cost, not
accuracy. The adversarial result is the one that survived every attempt
to find a cheaper substitute, and it is the claim we would defend if
forced to keep only one.

\paragraph{The real-task adversarial result is scoped to one repository.}
The pre-registered design was six paired worlds and we ran all six. The
burial magnitude replicates on \texttt{django} and is absent on
\texttt{sympy}, and the ordering that held in five worlds reverses in the
sixth, leaving $5/6$ and $p = 0.0625$ where six worlds were run to reach
$0.031$ (\S\ref{sec:swebench-attack}). Two sequences chosen for having
enough tasks cannot tell us which is typical, so the honest scope of the
real-task adversarial leg is one repository, and the mechanism argument
of \S\ref{sec:adversary} rests on the synthetic suite plus a
\texttt{django}-shaped confirmation rather than on a replicated law.


\pointer{Nine further threats to validity---among them vocabulary
coupling, the retriever bound on the real-task leg, and consolidation as
a laundering surface---are stated unabridged in \S\ref{sec:threats}.}
\inplace{limitations-threats}
